
\section{Implementation of Multi-Suction Cup Gripper Optimization}

This chapter presents the optimization framework developed for a multi-suction cup gripper
consisting of three suction cups arranged in a triangular configuration. The framework supports both
evolutionary multi-objective optimization and Bayesian multi-objective optimization. The
evolutionary algorithms are implemented using the \texttt{pymoo} library~\cite{pymoo}, whereas
Bayesian optimization is implemented using the \texttt{Ax} platform and its underlying
\texttt{BoTorch} modeling framework~\cite{olson2025ax,botorch}.

For both optimization approaches, the evaluation of candidate gripper layouts is delegated to an
external service through an API. This separation provides a modular architecture in which the
optimization strategy and the domain-specific grasp evaluation remain decoupled. Consequently, the
same decision-variable representation, grasp-generation procedure, and objective functions can be
used by both evolutionary and Bayesian optimization methods, allowing their performance to be
compared under a common evaluation pipeline.

\subsection{Reduced Parameterization of Multi-Suction Cup Gripper}

The gripper consists of three suction cups whose positions must ultimately be represented in
Cartesian space in order to perform grasp generation, collision checking, and objective computation.
However, the Cartesian suction cup coordinates are not used directly as optimization variables.
Instead, the layout is parameterized in a reduced form to constrain the search space and encode
geometric structure.

All three suction cups are constrained to lie on a common circle in the gripper plane. Their
configuration is therefore defined by one shared radius $r$ and three angular parameters
$\alpha_1,\alpha_2,\alpha_3$:
\begin{equation}
    \mathbf{q} = [r,\alpha_1,\alpha_2,\alpha_3]
\end{equation}
\nomenclature{$\mathbf{q}$}{Reduced decision variable vector defining the suction cup layout}
\nomenclature{$r$}{Common radial distance of the suction cups from the gripper center}
\nomenclature{$\alpha_i$}{Angular position of the $i$-th suction cup on the circle}

The angular parameters determine the positions of the suction cups on the circle, while the common
radius determines their distance from the gripper center. For a fixed gripper height $z_0$, the
Cartesian coordinates of the suction cups are obtained through the polar-to-Cartesian transformation
\begin{equation}
    \mathbf{c}_i =
    \begin{bmatrix}
        r\cos(\alpha_i) \\
        r\sin(\alpha_i) \\
        z_0
    \end{bmatrix},
    \qquad i=1,2,3.
\end{equation}
\nomenclature{$\mathbf{c}_i$}{Cartesian position of the $i$-th suction cup in the gripper frame}
\nomenclature{$z_0$}{Fixed height coordinate of the suction cups}

The resulting Cartesian gripper configuration is written as
\begin{equation}
    \mathbf{p} = [\mathbf{c}_1,\mathbf{c}_2,\mathbf{c}_3].
\end{equation}
\nomenclature{$\mathbf{p}$}{Cartesian suction cup configuration derived from the reduced
parameterization}

Thus, the optimization is carried out in the reduced 4-dimensional parameter space $\mathbf{q}$,
while the grasp evaluation operates on the derived Cartesian coordinates $\mathbf{p}$.

This formulation reduces the dimensionality of the search space from six planar degrees of freedom,
\begin{equation}
    [(x_1,y_1),(x_2,y_2),(x_3,y_3)] \in \mathbb{R}^6,
\end{equation}
to four variables,
\begin{equation}
    [r,\alpha_1,\alpha_2,\alpha_3].
\end{equation}
The reduction improves search efficiency and excludes many geometrically redundant or undesirable
layouts while preserving sufficient flexibility for diverse triangular suction cup arrangements.

\paragraph{Assumption}
A triangular configuration of three suction cups is assumed to yield a stable
grasp~\cite{madan2022realrobotchallenge2021}. This is motivated by the fact that three spatially
distributed contact points form a statically stable support polygon, enabling effective load
distribution and resistance to external disturbances.

\subsection{Decision Variables and Bounds for Multi-Suction Cup Gripper}

The optimization is performed over the reduced decision vector
\begin{equation}
    \mathbf{q} = [r,\alpha_1,\alpha_2,\alpha_3].
\end{equation}

All variables are treated as continuous. The angular variables are bounded by
\begin{equation}
    \alpha_i \in [0,2\pi), \qquad i=1,2,3,
\end{equation}
and the common radius is bounded by
\begin{equation}
    r \in [r_{\min},r_{\max}].
\end{equation}
\nomenclature{$r_{\min}$}{Lower bound on the common suction cup radius}
\nomenclature{$r_{\max}$}{Upper bound on the common suction cup radius}

The upper bound $r_{\max}$ is determined from the bounding box of the largest workpiece considered
during optimization, ensuring that the gripper footprint remains compatible with the object
dimensions. The lower bound $r_{\min}$ is determined from the diameter of an individual suction cup,
ensuring that the configuration remains physically realizable and that the suction cups do not
collapse into an excessively compact arrangement.

All optimization methods generate candidate values within these bounds. In the evolutionary
algorithms, the initial population is obtained by random sampling and subsequently modified through
the operators of the selected algorithm. In Bayesian optimization, Ax first evaluates the center of
the search space and then uses Sobol sampling to generate an initial space-filling design. After
sufficient initialization data have been collected, Ax transitions to model-based Bayesian
optimization.

\subsection{Optimization Framework of Multi-Suction Cup Gripper}

The multi-suction cup gripper problem is solved using either an evolutionary multi-objective
optimization algorithm or Bayesian multi-objective optimization. Both approaches operate on the same
reduced decision vector, variable bounds, objective functions, and external evaluation service.

The common optimization configuration contains:
\begin{itemize}
    \item Number of objectives $n_{\text{obj}}$
    \item Decision-variable definitions and bounds
    \item Optimization direction
    \item Choice of optimization method
    \item Evaluation endpoint URL
\end{itemize}

For evolutionary optimization, the additional configuration parameters include:
\begin{itemize}
    \item Population size $N$
    \item Maximum number of generations $G_{\max}$
    \item Choice of evolutionary algorithm
    \item Algorithm-specific parameters, such as crossover and mutation settings
\end{itemize}

For Bayesian optimization, the additional configuration parameters include:
\begin{itemize}
    \item Number of optimization batches $B_{\max}$
    \item Maximum number of candidates per batch $b$
    \item Bayesian generation strategy
\end{itemize}

In the current Bayesian implementation, the configuration value named \texttt{generations}
represents the number of candidate-generation batches rather than evolutionary generations. The
maximum requested number of Bayesian evaluations is therefore
\begin{equation}
    E_{\max}^{\mathrm{BO}} = B_{\max}b.
\end{equation}
The actual number may be smaller if Ax returns fewer than $b$ candidates in a particular batch.
\nomenclature{$B_{\max}$}{Maximum number of Bayesian optimization batches}
\nomenclature{$b$}{Maximum number of Bayesian candidates requested per batch}
\nomenclature{$E_{\max}^{\mathrm{BO}}$}{Maximum requested number of Bayesian objective evaluations}

\subsubsection{Supported Optimization Methods}

The framework supports the following evolutionary and Bayesian multi-objective optimization methods:

\begin{table}[htbp]
\centering
\caption{Supported multi-objective optimization methods.}
\label{tab:algorithms_sc}
\begin{tabularx}{\textwidth}{@{}l l >{\raggedright\arraybackslash}X@{}}
\toprule
\textbf{Family} & \textbf{Identifier} & \textbf{Method} \\
\midrule
Evolutionary & NSGA2 & Non-dominated Sorting Genetic Algorithm II~\cite{nsga2} \\
Evolutionary & NSGA3 & Non-dominated Sorting Genetic Algorithm III~\cite{nsga3} \\
Evolutionary & MOEAD & Multi-Objective Evolutionary Algorithm based on Decomposition~\cite{moead} \\
Evolutionary & CMOPSO & Competitive Multi-Objective Particle Swarm Optimization \\
Evolutionary & SMSEMOA & S-Metric Selection Evolutionary Multi-Objective Algorithm~\cite{smsemoa} \\
Bayesian & AX-MBM & Ax modular BoTorch model for multi-objective Bayesian
    optimization~\cite{olson2025ax, botorch} \\
\bottomrule
\end{tabularx}
\end{table}

For evolutionary algorithms requiring reference directions, such as NSGA-III and MOEA/D, uniform
reference directions are generated using the method of Das and Dennis~\cite{das_dennis}:
\begin{equation}
    \Lambda =
    \texttt{get\_reference\_directions}
    (\text{``das-dennis''},\; n_{\text{obj}},\; n_{\text{partitions}}).
\end{equation}
\nomenclature{$\Lambda$}{Set of reference directions for decomposition-based algorithms}
\nomenclature{$n_{\text{obj}}$}{Number of objectives}

The Bayesian implementation uses the default Ax generation strategy identified at runtime as
\texttt{Center+Sobol+MBM:fast}. This strategy consists of three stages:
\begin{enumerate}
    \item Evaluation of the center of the bounded search space
    \item Sobol quasi-random sampling for initialization
    \item Model-based optimization using the Ax modular BoTorch model (MBM)
\end{enumerate}

The Center and Sobol stages do not employ a surrogate model or acquisition function. After
transitioning to the MBM stage, Ax fits probabilistic surrogate models to the completed observations
and selects new candidates using an acquisition function appropriate for the multi-objective
problem. For continuous multi-objective problems, this generally involves Gaussian-process
surrogates and a hypervolume-improvement acquisition function. The exact model and
acquisition-function classes are selected by Ax according to the installed version and experiment
configuration and should therefore be recorded from the runtime configuration to ensure
reproducibility.

\subsection{Initialization of Multi-Suction Cup Gripper Optimization}

The initialization procedure depends on the selected optimization approach.

\subsubsection{Evolutionary Initialization}

For evolutionary optimization, an initial population of $N$ candidate solutions is generated by
sampling the reduced decision variables within their bounds:
\begin{equation}
    \mathcal{P}^{(0)}
    =
    \{\mathbf{q}_1,\mathbf{q}_2,\dots,\mathbf{q}_N\}.
\end{equation}
\nomenclature{$\mathcal{P}^{(0)}$}{Initial population of candidate decision vectors}
\nomenclature{$N$}{Population size}

Each candidate decision vector $\mathbf{q}_k$ is transformed into its corresponding Cartesian
gripper configuration $\mathbf{p}_k$ before objective evaluation.

\subsubsection{Bayesian Initialization}

Bayesian optimization does not maintain an evolutionary population. Instead, it stores a dataset of
completed trials:
\begin{equation}
    \mathcal{D}_t =
    \left\{
        \left(\mathbf{q}_k,\mathbf{g}(\mathbf{q}_k)\right)
    \right\}_{k=1}^{t},
\end{equation}
where $\mathbf{q}_k$ is an evaluated gripper parameterization and $\mathbf{g}(\mathbf{q}_k)$ is its
observed objective vector. \nomenclature{$\mathcal{D}_t$}{Dataset of completed Bayesian optimization
trials after $t$ evaluations}

The Ax generation strategy first evaluates the center of the search space and then generates an
initial set of candidates using Sobol quasi-random sampling. Sobol sampling provides broad coverage
of the bounded decision space before a surrogate model is fitted. Once the initialization
requirements of the generation strategy have been satisfied, Ax transitions to model-based candidate
generation.

\subsection{Objective Function Evaluation of Multi-Suction Cup Gripper via API}

The evaluation of objective values is decoupled from the optimization loop and delegated to an
external evaluation service. For each candidate solution $\mathbf{q}$, the corresponding Cartesian
suction cup coordinates are first computed by transforming the radial and angular parameters into 3D
positions. The resulting Cartesian gripper configuration is then passed to the evaluation endpoint,
which executes the grasp evaluation procedure and returns the objective values.

This process can be written as
\begin{equation}
    \mathbf{p} = T(\mathbf{q}),
    \qquad
    \mathbf{f}(\mathbf{q}) = \mathbf{f}(T(\mathbf{q}))
    = \texttt{API\_call}(\mathbf{p})
    \rightarrow
    [f_1(\mathbf{p}),\; f_2(\mathbf{p})].
\end{equation}
where $T(\cdot)$ denotes the transformation from the reduced parameterization to Cartesian suction
cup coordinates. \nomenclature{$T(\mathbf{q})$}{Transformation from reduced decision variables to
Cartesian suction cup coordinates}

In the implementation, this transformation is carried out programmatically in the evaluation
pipeline before grasp generation and objective computation. For a candidate radius and list of
angles, the suction cup coordinates are constructed as
\begin{equation}
    \mathbf{c}_i = (r\cos(\alpha_i),\; r\sin(\alpha_i),\; z_0),
    \qquad i=1,2,3.
\end{equation}

This architecture enables:
\begin{itemize}
    \item Separation of concerns between optimization logic and domain-specific evaluation
    \item Scalability through independent deployment of the evaluation service
    \item Reusability of the optimization framework for different problem domains
\end{itemize}

\subsubsection{Grasp Evaluation}
\label{sec:grasp_eval_sc}

The evaluation service receives a Cartesian gripper configuration $\mathbf{p} =
[\mathbf{c}_1,\mathbf{c}_2,\mathbf{c}_3]$ and computes the associated grasp quality measures through
the following procedure.

\begin{enumerate}
    \item \textbf{Candidate Grasp Generation and Feasibility Filtering}

    Suction grasp candidates are generated using a method developed by Fraunhofer IPA, following the
    automatic suction grasp generation procedure described by ~\cite{Kleeberger}. Points are sampled
    on the surface of the workpiece, and each sampled point defines a candidate grasp through its
    position and local surface normal. The local surface normal determines the approach direction of
    the suction cup.

    In the original method, each suction grasp candidate is tested for feasible seal formation using
    a quasi-static spring model of a suction cup, as introduced in ~\cite{SC_quasi_static}. If the
    grasp point is considered feasible, the candidate is further tested with a gripper collider
    model for collisions between the gripper and the object.

    In this thesis, the method was extended to a gripper with three suction cups. For each candidate
    gripper pose, the seal-formation test is performed for all three suction cup contact positions
    $\mathbf{c}_1$, $\mathbf{c}_2$, and $\mathbf{c}_3$. A candidate is considered feasible only if
    all three suction cups form a sufficient seal. The candidate is then checked for collisions in
    simulation, and only collision-free candidates are retained as valid grasps.
    \nomenclature{$\mathbf{p}$}{Cartesian gripper configuration defined by the three suction cup
    contact positions} \nomenclature{$\mathbf{c}_i$}{Cartesian contact position of suction cup $i$,
    with $i \in \{1,2,3\}$}

    \item \textbf{Clustering of Grasps}

    The set $\mathcal{G}_{\mathrm{valid}}$ often contains many grasps that are nearly identical in
    position and orientation. When the number of valid grasps $m = |\mathcal{G}_{\mathrm{valid}}|$
    exceeds a threshold $\texttt{nb\_grasps\_clustering}$, a farthest-point clustering procedure is
    applied to select a diverse reduced subset. Otherwise, all valid grasps are used directly.

    Each grasp $g_i$ is represented by a pose descriptor
    \begin{equation}
        \mathbf{z}_i =
        [\,\mathrm{vec}(\mathbf{R}_i)^\top,\ \mathbf{t}_i^\top\,]^\top
        \in \mathbb{R}^{12},
    \end{equation}
    where $\mathbf{R}_i \in SO(3)$ is the orientation and $\mathbf{t}_i \in \mathbb{R}^3$ is the
    position. Two distances are used: the translational distance
    \begin{equation}
        d_{\mathrm{pos}}(i,j) = \|\mathbf{t}_i - \mathbf{t}_j\|_2
    \end{equation}
    and the orientation distance
    \begin{equation}
        d_{\mathrm{ori}}(i,j) =
        \frac{1}{9}
        \|\mathrm{vec}(\mathbf{R}_i) - \mathrm{vec}(\mathbf{R}_j)\|_2^2.
    \end{equation}

    Cluster centers are selected greedily by farthest-point sampling in position space. The first
    center is the grasp with the largest mean translational distance to all others, and each
    subsequent center maximizes its minimum distance to the already selected centers. This promotes
    broad spatial coverage.

    Optionally, orientation diversity is added within each center's local neighborhood. For a
    selected center $m_k$, grasps within a spatial radius $\delta_{\mathrm{intra}}$ form a
    neighborhood, from which up to $q_{\mathrm{intra}}$ additional grasps are chosen by
    farthest-point sampling on $d_{\mathrm{ori}}$, capturing distinct orientations at similar
    positions.

    The resulting clustered subset
    \begin{equation}
        \mathcal{G}_{\mathrm{clustered}} \subseteq \mathcal{G}_{\mathrm{valid}}
    \end{equation}
    is then used for diversity evaluation. \nomenclature{$\mathcal{G}_{\mathrm{clustered}}$}{Reduced
    clustered grasp set used for diversity evaluation}
    \nomenclature{$d_{\mathrm{pos}}(i,j)$}{Pairwise translational distance between grasp poses $i$
    and $j$} \nomenclature{$d_{\mathrm{ori}}(i,j)$}{Pairwise orientation distance between grasp
    poses $i$ and $j$} \nomenclature{$\delta_{\mathrm{intra}}$}{Spatial threshold for selecting
    additional local representatives} \nomenclature{$q_{\mathrm{intra}}$}{Maximum number of local
    orientation-diverse representatives}

    \item \textbf{Triangle Area Calculation}

    The area of the triangle formed by the three suction cups is computed from their Cartesian
    coordinates. If the planar coordinates of the three cups are $(x_1,y_1)$, $(x_2,y_2)$, and
    $(x_3,y_3)$, then
    \begin{equation}
        A(\mathbf{p}) =
        \frac{1}{2}
        \left|
            (x_2 - x_1)(y_3 - y_1) - (x_3 - x_1)(y_2 - y_1)
        \right|.
    \end{equation}

    This objective favors gripper layouts with a wider support polygon, which is associated with
    improved grasp stability. \nomenclature{$A(\mathbf{p})$}{Area of the triangle formed by the
    three suction cups}


\item \textbf{Diversity Metric Calculation}

A diversity score $D(\mathbf{p})$ is computed to quantify how broadly the resulting valid grasp set
is distributed in position and orientation.

Let the set used for diversity evaluation contain $n$ grasps. If clustering is performed, then
\begin{equation}
    n = |\mathcal{G}_{\mathrm{clustered}}|,
\end{equation}
otherwise
\begin{equation}
    n = |\mathcal{G}_{\mathrm{valid}}|.
\end{equation}

Let the corresponding grasp positions be
\begin{equation}
    \{\mathbf{y}_1,\dots,\mathbf{y}_n\}
\end{equation}
and the corresponding orientation matrices be
\begin{equation}
    \{\mathbf{R}_1,\dots,\mathbf{R}_n\}.
\end{equation}

To normalize the spatial term with respect to object size, a characteristic object scale is defined
as the Euclidean norm of the axis-aligned bounding-box extents:
\begin{equation}
    s = \|\mathbf{e}\|_2,
\end{equation}
where $\mathbf{e} \in \mathbb{R}^3$ contains the bounding-box side lengths.
\nomenclature{$s$}{Characteristic object scale derived from the mesh bounding box}
\nomenclature{$\mathbf{e}$}{Axis-aligned bounding-box side lengths of the object}
\nomenclature{$\mathbf{y}_i$}{Position of grasp $i$ used in diversity evaluation}

The spatial component of diversity is computed as the mean pairwise Euclidean distance between grasp
positions:
\begin{equation}
    D_{\mathrm{pos}}(\mathbf{p}) =
    \frac{1}{s}
    \cdot
    \frac{1}{n(n-1)}
    \sum_{\substack{i,j=1 \\ i \neq j}}^{n}
    \|\mathbf{y}_i - \mathbf{y}_j\|_2.
\end{equation}

The orientation component is computed from the mean pairwise geodesic distance between rotation
matrices. For two grasp orientations $\mathbf{R}_a,\mathbf{R}_b \in SO(3)$, the relative rotation is
\begin{equation}
    \Delta \mathbf{R} = \mathbf{R}_a^\top \mathbf{R}_b
\end{equation}
and the associated geodesic distance is
\begin{equation}
    d_{\mathrm{geo}}(\mathbf{R}_a,\mathbf{R}_b)
    =
    \cos^{-1}\left(
        \frac{\mathrm{tr}(\Delta \mathbf{R}) - 1}{2}
    \right).
\end{equation}

For numerical stability, the argument of $\cos^{-1}$ is clipped to the interval $[-1,1]$.

The normalized orientation diversity is then defined as
\begin{equation}
    D_{\mathrm{rot}}(\mathbf{p}) =
    \frac{1}{\pi}
    \cdot
    \frac{2}{n(n-1)}
    \sum_{1 \leq i < j \leq n}
    d_{\mathrm{geo}}(\mathbf{R}_i,\mathbf{R}_j).
\end{equation}

The final diversity score is computed as the equally weighted combination
\begin{equation}
    D(\mathbf{p}) =
    \mathrm{clip}\left(
        0.5\,D_{\mathrm{pos}}(\mathbf{p}) +
        0.5\,D_{\mathrm{rot}}(\mathbf{p}),
        0, 1
    \right).
\end{equation}
\nomenclature{$D(\mathbf{p})$}{Diversity metric of the grasp distribution}
\nomenclature{$D_{\mathrm{pos}}(\mathbf{p})$}{Spatial component of the diversity metric}
\nomenclature{$D_{\mathrm{rot}}(\mathbf{p})$}{Orientation component of the diversity metric}
\nomenclature{$d_{\mathrm{geo}}(\mathbf{R}_a,\mathbf{R}_b)$}{Geodesic distance between two grasp
orientations} \nomenclature{$\Delta\mathbf{R}$}{Relative rotation between two grasp orientations}
\nomenclature{$n$}{Number of grasps used for diversity evaluation}

\paragraph{Continuity considerations}
Clustering is only applied when the number of valid grasps exceeds the predefined threshold
$\texttt{nb\_grasps\_clustering}$. If fewer valid grasps are available, clustering is skipped and
the diversity metric is evaluated directly on the valid grasp set. Therefore, the omission of
clustering does not by itself introduce a discontinuity in the objective function.

A discontinuity may arise only in the degenerate case where fewer than two valid grasps remain after
collision filtering. Since both the spatial and orientation components of the diversity metric are
based on pairwise distances, these terms cannot be evaluated meaningfully for $n<2$. In that case,
the corresponding terms are set to zero. As a result, small changes in the gripper configuration may
cause an abrupt transition between a nonzero diversity score and a degenerate zero-valued diversity
contribution.

To reduce the likelihood of this case, the number of initially sampled surface points was increased
from 100 to 1000. The denser surface sampling produces a larger set of candidate grasps and makes it
much less likely that fewer than two valid grasps remain after filtering. Consequently, the
diversity metric is evaluated on a sufficiently rich grasp set in the vast majority of cases,
yielding a more stable optimization objective in practice.
\end{enumerate}

\subsubsection{Objective Evaluation}

After grasp evaluation, each candidate layout is assigned the natural objective vector
\begin{equation}
    \mathbf{g}(\mathbf{q})
    =
    [A(\mathbf{p}),\;D(\mathbf{p})],
    \qquad
    \mathbf{p}=T(\mathbf{q}).
\end{equation}
\nomenclature{$\mathbf{g}(\mathbf{q})$}{Natural multi-objective function vector containing triangle
area and grasp diversity}

Both components are maximized. A larger triangle area is associated with a wider support polygon and
improved grasp stability, while a larger diversity score indicates that the valid grasps are more
broadly distributed in position and orientation.

The objective-direction convention is handled differently by the two optimization frameworks. Since
\texttt{pymoo} assumes minimization, the evolutionary algorithms negates objective
vector
\begin{equation}
    \mathbf{f}_{\mathrm{EA}}(\mathbf{q})
    =
    [-A(\mathbf{p}),\;-D(\mathbf{p})].
\end{equation}
\nomenclature{$\mathbf{f}_{\mathrm{EA}}(\mathbf{q})$}{Minimization-form objective vector supplied to
the evolutionary algorithms}

In the Bayesian implementation, the external evaluation service values are left as is, since
\texttt{Ax} assumes maximization
\begin{equation}
    \mathbf{f}_{\mathrm{BO}}(\mathbf{q})
    =
    [A(\mathbf{p}),\;D(\mathbf{p})].
\end{equation}
\nomenclature{$\mathbf{f}_{\mathrm{BO}}(\mathbf{q})$}{Objective observations supplied to Ax}

The optimization direction is specified through the Ax objective expression. Consequently, the
observations themselves are not negated before being reported to Ax. This preserves the original
physical meaning of the objective values in the returned Pareto set.


\subsection{Optimization Loop for Multi-Suction Cup Gripper}

The candidate-generation process differs between evolutionary and Bayesian optimization, while the
external candidate-evaluation procedure remains identical.

\subsubsection{Evolutionary Optimization Loop}

The evolutionary optimization proceeds according to
\begin{equation}
    \mathcal{P}^{(t+1)}
    =
    \mathcal{A}
    \left(
        \mathcal{P}^{(t)},
        \mathbf{f}_{\mathrm{EA}}
    \right),
    \qquad
    t=0,1,\dots,
\end{equation}
where $\mathcal{A}$ denotes the selected evolutionary algorithm.
\nomenclature{$\mathcal{A}$}{Selected multi-objective evolutionary algorithm}
\nomenclature{$t$}{Evolutionary generation index}

Depending on the selected algorithm, the new population is generated through operations such as
selection, crossover, mutation, decomposition, particle movement, or survival based on hypervolume
contribution. Every candidate in the resulting population is evaluated through the common API-based
grasp evaluation pipeline.

\subsubsection{Bayesian Optimization Loop}

The Bayesian optimization procedure maintains the observation dataset $\mathcal{D}_t$. During the
initialization phase, candidates are generated from the center of the search space and by Sobol
sampling. During the model-based phase, the available observations are used to fit probabilistic
surrogate models:
\begin{equation}
    \widehat{\mathbf{g}}_t
    =
    \mathcal{M}(\mathcal{D}_t),
\end{equation}
where $\mathcal{M}$ denotes the surrogate-modeling procedure used by Ax and BoTorch.
\nomenclature{$\mathcal{M}$}{Bayesian surrogate-modeling procedure}
\nomenclature{$\widehat{\mathbf{g}}_t$}{Probabilistic surrogate approximation after $t$ evaluations}

The surrogate models provide predictions of the objective values together with their uncertainty. An
acquisition function $\alpha(\mathbf{q}\mid\mathcal{D}_t)$ uses this information to balance the
exploration of uncertain regions and the exploitation of regions expected to improve the current
Pareto frontier.

For a requested batch size $b$, the next candidate batch is generated as
\begin{equation}
    \mathcal{Q}_{t+1}
    =
    \operatorname*{arg\,max}_{
        \substack{
            \mathcal{Q}\subset\mathcal{X}\\
            |\mathcal{Q}|\leq b
        }
    }
    \alpha(\mathcal{Q}\mid\mathcal{D}_t),
\end{equation}
where $\mathcal{X}$ is the bounded decision space. \nomenclature{$\mathcal{Q}_{t+1}$}{Batch of
candidates generated by Bayesian optimization} \nomenclature{$\mathcal{X}$}{Bounded decision space}
\nomenclature{$\alpha$}{Bayesian acquisition function}

Each candidate $\mathbf{q}\in\mathcal{Q}_{t+1}$ is evaluated through the external evaluation
service. The resulting observation is then supplied to Ax as a completed trial:
\begin{equation}
    \mathcal{D}_{t+1}
    =
    \mathcal{D}_t
    \cup
    \left\{
        \left(
            \mathbf{q}_{t+1},
            \mathbf{g}(\mathbf{q}_{t+1})
        \right)
    \right\}.
\end{equation}

The surrogate model is subsequently updated using the expanded dataset, and the process is repeated
until the specified evaluation budget is exhausted. Although multiple candidates can be requested in
one batch, the current implementation submits the candidates in that batch sequentially to the
evaluation service.

\subsection{Termination Criteria for Multi-Suction Cup Gripper Optimization}

The termination criteria depend on the selected optimization approach.

\begin{table}[htbp]
\centering
\caption{Termination criteria for evolutionary and Bayesian optimization.}
\label{tab:termination_sc}
\begin{tabularx}{\textwidth}{@{}l l l >{\raggedright\arraybackslash}X@{}}
\toprule
\textbf{Approach} & \textbf{Criterion} & \textbf{Symbol} & \textbf{Description} \\
\midrule
Evolutionary & Design-space tolerance & $x_{\text{tol}}$ & Change in decision variables below
threshold \\

Evolutionary & Objective-space tolerance & $f_{\text{tol}}$ & Relative objective improvement below
threshold \\

Evolutionary & Maximum generations & $G_{\max}$ & Maximum number of generations reached \\

Evolutionary & Maximum evaluations & $E_{\max}$ & Maximum number of function evaluations reached \\

Evolutionary & Hypervolume & $HV_{\max}$ & Hypervolume progression stagnates  \\

Bayesian & Maximum batches & $B_{\max}$ & Maximum number of candidate-generation batches reached \\

Bayesian & Maximum evaluations & $E_{\max}^{\mathrm{BO}}$ & Bayesian evaluation budget exhausted \\
\bottomrule
\end{tabularx}
\end{table}

For evolutionary optimization, the design-space and objective-space tolerances are evaluated over a
sliding window of $w$ generations. \nomenclature{$x_{\text{tol}}$}{Design-space convergence
tolerance} \nomenclature{$f_{\text{tol}}$}{Objective-space convergence tolerance}
\nomenclature{$G_{\max}$}{Maximum number of evolutionary generations}
\nomenclature{$E_{\max}$}{Maximum number of evolutionary objective evaluations}

In the current Bayesian implementation, termination is budget-based. The optimization performs at
most $B_{\max}$ candidate-generation batches, with up to $b$ candidates requested per batch.
Therefore,
\begin{equation}
    E_{\max}^{\mathrm{BO}} = B_{\max}b.
\end{equation}
No design-space or objective-space convergence tolerance is currently applied to the Bayesian
optimization loop.

\subsection{Performance Monitoring for Multi-Suction Cup Gripper}

The hypervolume indicator~\cite{hypervolume} is tracked at each generation via a callback mechanism:
\begin{equation}
    \text{HV}^{(t)} = \text{HV}(\mathcal{F}^{(t)},\mathbf{r})
\end{equation}
\nomenclature{$\text{HV}^{(t)}$}{Hypervolume indicator at generation $t$}
\nomenclature{$\mathbf{r}$}{Reference point for hypervolume computation}
\nomenclature{$\mathcal{F}^{(t)}$}{Objective values of the population at generation $t$}

Here, $\mathbf{r} \in \mathbb{R}^{n_{\text{obj}}}$ is the reference point and $\mathcal{F}^{(t)}$
denotes the set of objective vectors at generation $t$. A monotonically increasing hypervolume
indicates convergence toward the Pareto front.

\subsection{Optional Hyperparameter Optimization for Multi-Suction Cup Gripper Problem}

The Evolutionary Optimization framework optionally supports automated tuning of algorithm
hyperparameters, such as crossover probability and mutation rate, using a meta-optimization
approach. A \texttt{HyperparameterProblem} is constructed around the selected algorithm and
evaluated over multiple random seeds. The mean hypervolume across runs serves as the meta-objective:
\begin{equation}
    \min_{\boldsymbol{\theta}}
    \left\{
        -\frac{1}{|\mathcal{Z}|}
        \sum_{s \in \mathcal{Z}} \text{HV}(\mathcal{F}_s,\mathbf{r})
    \right\}
\end{equation}
\nomenclature{$\boldsymbol{\theta}$}{Algorithm hyperparameters} \nomenclature{$\mathcal{Z}$}{Set of
random seeds for meta-optimization}

The best hyperparameters are then applied to the main optimization run.

\subsection{Output of Multi-Suction Cup Gripper Optimization}

The optimization produces a set of Pareto-optimal solutions:
\begin{equation}
    \mathcal{P}^* =
    \left\{
        \mathbf{q}
        \;\middle|\;
        \nexists\; \mathbf{q}' \in \mathcal{P} :
        \mathbf{f}(\mathbf{q}') \preceq \mathbf{f}(\mathbf{q})
    \right\}.
\end{equation}
\nomenclature{$\mathcal{P}^*$}{Set of Pareto-optimal solutions in the reduced decision space}

The final output includes:
\begin{itemize}
    \item The Pareto-optimal set $\mathcal{P}^*$ in the reduced decision space, with corresponding
    objective values
    \item The associated radial gripper parameterizations $\mathbf{q} =
    [r,\alpha_1,\alpha_2,\alpha_3]$
    \item The hypervolume progression $\{\text{HV}^{(0)},\text{HV}^{(1)},\dots,\text{HV}^{(T)}\}$
    \item Total optimization wall-clock time
\end{itemize}

These solutions represent optimal trade-offs between grasp diversity and triangle area, i.e., grasp
robustness and stability.